\color{unimain}{\subsection{Логика отключения}}
\color{black}

\begin{enumerate}

\item
Функциональный блок логики отключения предназначен для формирования общего сигнала отключения при срабатывании функций РЗА и общего сигнала запрета АПВ.  

\item
Функциональная схема общей логики отключения пприведена на рисунке \ref{pict_lo:logic_otkl_obsh}.

\vspace{2mm}
\begin{figure}[H]
\centering
\includegraphics[width=0.75\textwidth,height=0.75\textheight,keepaspectratio]{\fbpath/03.01 ЛВ/761.4011 ЛО/img/logic_otkl.pdf}
\captionof{figure}{Функциональная схема общей логики отключения}\label{pict_lo:logic_otkl_obsh}
\end{figure}

\item
Команда логики отключения при срабатывании функций РЗА, действующая на отключение выключателей \textit{«Ср. защит»} формируется:
\begin{itemize}
    \item при срабатывании внутренних защит устройства (сигналы \textit{«Срабатывание»}, \textit{«ТУ ДЗ: Передача»}, \textit{«Срабатывание ЛОСК»});
    \item при наличии сигнала срабатывания внешних защит (сигнал \textit{«Ср. внеш. защит»});
    \item при наличии сигнала срабатывания основных защит (сигнал \textit{«Ср. осн. защит»});
    \item при наличии внешних сигналов УРОВ (сигнал \textit{«Внеш. УРОВ от РЗА»}).
\end{itemize}

\item
Команда останова ВЧ приемопередатчика \textit{«Останов ВЧ»} формируется:
\begin{itemize}
    \item при формировании сигнала логики отключения при срабатывании функций РЗА (сигналы \textit{«Ср. защит»});
    \item при наличии сигналов срабатывания УРОВ (сигнал \textit{«УРОВ: Срабатывание»}).
\end{itemize}

\item
Команда телеускорения отключения трёх фаз (ТУ ОТФ) \textit{«Ср. защит на ТУ ОТФ»} формируется при наличии сигнала логики отключения при срабатывании функций РЗА \textit{«Ср. защит»} и отсутствии сигнала срабатывания токовой защиты ошиновки \textit{«ТЗО: Срабатывание»}.

\item
Формирование выходного сигнала пуска защит \textit{«Пуск защит»} осуществляется при наличии внешнего сигнала \textit{«П. осн. защиты»}, либо при пуске ДЗ, ТЗНП, МФТО или АМТЗ.

\item
При наличии входного сигнала вывода логики отключения \textit{«ЛО: Выведено»} блокируется дальнейшее прохождение сигналов для всех выходных сигналов данного функционального блока.

\item
Функциональная схема общей логики запрета АПВ приведена на рисунке \ref{pict_lo:logic_zapv}. Уставки логики запрета АПВ в таблице \ref{app_set:lo}.

\vspace{2mm}
\begin{figure}[H]
\centering
\includegraphics[width=0.85\textwidth,height=0.85\textheight,keepaspectratio]{\fbpath/03.01 ЛВ/761.4011 ЛО/img/zapv.pdf}
\captionof{figure}{Функциональная схема общей логики запрета АПВ}\label{pict_lo:logic_zapv}
\end{figure}

\item
Команда на запрет АПВ \textit{«Общий запрет АПВ»} формируется:
\begin{itemize}
	\item при срабатывании ступеней внутренних защит устройства (при введеном соответствующем программном переключателе);
    \item при срабатывании УРОВ;
    \item при приеме внешнего сигнала запрета АПВ;
    \item при местном управлении.
\end{itemize}

\end{enumerate}